%% parts/p03/types-ultraproducts.tex --- P3 types, saturation, and
%% ultraproducts.  Migrates v1.1 lines 6510--7149.  Statements remain in
%% the body; all retained source proofs are compressed in the parallel annex.
\chapter{Types, Saturation, and Ultraproducts}
\label{sec:p03-types-ultraproducts}

Throughout this chapter, structures are single-sorted first-order structures
with equality and nonempty domains.  Uses of the Axiom of Choice or of the
weaker Ultrafilter Lemma are stated at the result that needs them.  Standard
references for the chapter are \cite{ChangKeisler1990,Hodges1993,Marker2002};
the ultraproduct theorem originates with \L o\'s \cite{Los1955}.

\section{Types over parameters and theories}
\label{sec:p03-types}

\begin{hbdefinition}{def:p03-type}{Type over a parameter set}
Let \(\mathcal M\) be a \(\Lang\)-structure, let
\(A\subseteq\dom\mathcal M\), and use the parameter expansion
\(\mathcal M_A\) of Definition~\ref{def:parameter-expansion}.  For a
finite tuple \(\bar x=(x_1,\ldots,x_n)\) of distinct variables, an
\emph{\(n\)-type over \(A\) consistent with \(\mathcal M\)} is a set
\(p(\bar x)\) of \(\Lang(A)\)-formulas, with free variables among \(\bar x\),
such that
\[
 \operatorname{Th}_{\Lang(A)}(\mathcal M_A)
 \cup\{\varphi(\bar d):\varphi(\bar x)\in p\}
\]
is satisfiable, where \(\bar d\) is a new \(n\)-tuple of constants.  It is
\emph{complete} when, for every such \(\varphi\), exactly one of
\(\varphi\) and \(\neg\varphi\) belongs to \(p\).  Thus ``consistent with''
names satisfiability in the expanded language, equivalently derivability
consistency by first-order soundness and completeness; it does not mean
realization in \(\mathcal M\).
\index{type!over a parameter set}
\index{complete type}
\end{hbdefinition}

\begin{hbdefinition}{def:p03-type-space}{Types over a theory and the Stone space}
Let \(T\) be a satisfiable \(\Lang\)-theory.  An \emph{\(n\)-type consistent
with \(T\)} is a set \(p(\bar x)\) of \(\Lang\)-formulas such that
\(T\cup\{\varphi(\bar d):\varphi\in p\}\) is satisfiable for new constants
\(\bar d\).  Let \(S_n(T)\) be the set of complete such types.  For each
formula \(\varphi(\bar x)\), put
\[
 [\varphi]=\{p\in S_n(T):\varphi\in p\}.
\]
The sets \([\varphi]\) form a clopen basis, since
\([\varphi]\cap[\psi]=[\varphi\land\psi]\) and
\(S_n(T)\setminus[\varphi]=[\neg\varphi]\).  This topology makes
\(S_n(T)\) the \emph{\(n\)-type space of \(T\)}.  Its points are complete
descriptions consistent with \(T\), not merely tuples in one selected model;
the distinction is needed for compactness, isolation, and omission below.
\index{type space}
\index{Stone space!of types}
\end{hbdefinition}

\begin{hbtheorem}{thm:p03-type-space-compact}{proved=annex:p3.types.1}{Compactness of type spaces}
Assume the Axiom of Choice in the ambient set theory.  If \(T\) is
satisfiable, then \(S_n(T)\) is compact and Hausdorff and has a basis of
clopen sets.
\end{hbtheorem}

For the complete theory of a two-element structure in the equality-only
language, \(S_1(T)\) has one point and \(S_2(T)\) has two: equality and
inequality.  Both points of \(S_2(T)\) are clopen, so that space is discrete.

\begin{hbdefinition}{def:p03-realization}{Realization and omission}
A tuple \(\bar b\in(\dom\mathcal M)^n\) \emph{realizes} a type
\(p(\bar x)\) over parameters from \(\mathcal M\) when
\(\mathcal M\entails\varphi(\bar b)\) for every \(\varphi\in p\).
The structure \(\mathcal M\) \emph{omits} \(p\) when no tuple in its domain
realizes \(p\).  Realization is therefore a relation between a tuple and a
type; finite satisfiability alone does not imply realization in the chosen
structure.
\index{type!realization}
\index{type!omission}
\end{hbdefinition}

For example, in the parameter expansion of \((\N,<)\), the type
\(\{c_n<x:n\in\N\}\) is finitely realized but not realized in \(\N\); under
Choice, compactness realizes it in an elementary extension.

\begin{hbdefinition}{def:p03-saturation}{Saturation}
Let \(\kappa\) be an infinite cardinal.  A structure \(\mathcal M\) is
\emph{\(\kappa\)-saturated} when, for every
\(A\subseteq\dom\mathcal M\) with \(\abs{A}<\kappa\), every complete type
in finitely many variables over \(A\) consistent with \(\mathcal M\) is
realized in \(\mathcal M\).  An infinite \(\mathcal M\) is
\emph{saturated} when it is \(\abs{\dom\mathcal M}\)-saturated.  The finite
tuple convention is part of this definition; sources allowing tuples of
length below \(\kappa\) use a stronger convention.
\index{saturation}
\index{kappa-saturation@\(\kappa\)-saturation}
\end{hbdefinition}

\begin{hbtheorem}{thm:p03-saturated-extension}{cited=Marker2002}{Saturated elementary extensions}
Assume the Axiom of Choice in the ambient set theory.  For every structure
\(\mathcal M\) and every infinite cardinal \(\kappa\), there is a
\(\kappa\)-saturated elementary extension
\(\mathcal N\succcurlyeq\mathcal M\).  No prescribed cardinal bound on
\(\mathcal N\) is asserted.
\end{hbtheorem}

The cited construction iterates elementary extensions that realize types over
small parameter sets and takes elementary-chain unions at limit stages; a
cardinality bound would require additional bookkeeping not supplied here.

\section{Homogeneity}
\label{sec:p03-homogeneity}

\begin{hbdefinition}{def:p03-homogeneity}{Homogeneity}
Let \(\kappa\) be an infinite cardinal.  A structure \(\mathcal M\) is
\emph{\(\kappa\)-homogeneous} when, whenever
\(A\subseteq\dom\mathcal M\) has \(\abs A<\kappa\), an elementary map
\(h:A\to\dom\mathcal M\) and an element \(b\in\dom\mathcal M\) admit an
elementary extension \(h':A\cup\{b\}\to\dom\mathcal M\).  Homogeneity is a
property of a structure extending partial elementary self-maps; unlike
saturation, it does not itself quantify over realization of every small type.
The next theorem is what connects the two notions.
\index{homogeneity}
\index{kappa-homogeneity@\(\kappa\)-homogeneity}
\end{hbdefinition}

\begin{hbtheorem}{thm:p03-saturation-homogeneity}{proved=annex:p3.types.2}{Saturation supplies homogeneity}
Assume the Axiom of Choice in the ambient set theory.  Every
\(\kappa\)-saturated structure is \(\kappa\)-homogeneous.
\end{hbtheorem}

Thus a saturated elementary extension not only realizes each small type: an
elementary map on fewer than \(\kappa\) elements can be extended by any one
further element.  This is the extension property used by back-and-forth.

\section{Principal types and omission}
\label{sec:p03-omitting-types}

\begin{hbdefinition}{def:p03-principal-type}{Principal, nonprincipal, and isolated types}
Let \(p(\bar x)\) be a type consistent with a theory \(T\).  It is
\emph{principal over \(T\)} when some formula \(\theta(\bar x)\) satisfies
both
\[
 T\cup\{\exists\bar x\,\theta(\bar x)\}\text{ is satisfiable},
 \qquad
 T\entails\forall\bar x\,(\theta\to\varphi)
 \quad(\varphi\in p).
\]
Otherwise it is \emph{nonprincipal over \(T\)}.  For a complete type this
is exactly the topological condition that its point of \(S_n(T)\) is isolated:
then \(\{p\}=[\theta]\).  Accordingly, \emph{isolated complete type} is the
topological name for the already defined principal complete type, not a
second classification.  This distinction from arbitrary consistent types is
the hypothesis used by the omitting-types theorem and by prime-model theory.
\index{principal type}
\index{nonprincipal type}
\index{isolated type}
\end{hbdefinition}

\begin{hbtheorem}{thm:p03-omitting-types}{cited=Hodges1993}{Omitting types theorem}
Let \(T\) be a satisfiable theory in a countable signature, and let
\(p_0,p_1,\ldots\) be a countable family of nonprincipal types consistent
with \(T\), each in a finite tuple of variables.  Then \(T\) has a countable
model omitting every \(p_n\).
\end{hbtheorem}

For the language with constants \(c_0,c_1,\ldots\), the theory saying that
they are pairwise distinct has the nonprincipal type
\(p(x)=\{x\ne c_n:n\in\N\}\).  The countable model consisting exactly of the
named elements omits \(p\), whereas an extension with an unnamed element
realizes it.  This separates nonprincipality of a type from omission by a
particular model.

\section{Filters and nonprincipal ultrafilters}
\label{sec:p03-ultrafilters}

\begin{hbdefinition}{def:p03-filter}{Filters and ultrafilters}
Let \(I\ne\varnothing\).  A \emph{filter} on \(I\) is a family
\(\mathcal F\subseteq\mathcal P(I)\) such that \(I\in\mathcal F\),
\(\varnothing\notin\mathcal F\), it is closed under binary intersections,
and it is upward closed within \(I\).  An \emph{ultrafilter} is a filter
\(\mathcal U\) such that, for every \(X\subseteq I\), exactly one of \(X\)
and \(I\setminus X\) lies in \(\mathcal U\).  It is \emph{principal at
\(i_0\)} when \(\mathcal U=\{X\subseteq I:i_0\in X\}\), and
\emph{nonprincipal} otherwise.  A subset of \(I\) is \emph{cofinite} when
its complement is finite.  These notions distinguish the large-index sets
used below from arbitrary families of coordinates.
The exactly-one condition is called the \emph{ultrafilter alternative}.
\index{filter}
\index{ultrafilter}
\index{ultrafilter alternative}
\index{ultrafilter!principal}
\index{ultrafilter!nonprincipal}
\end{hbdefinition}

\begin{hblemma}{thm:p03-nonprincipal-ultrafilter}{proved=annex:p3.types.3}{Existence and characterization of nonprincipal ultrafilters}
Assume the Ultrafilter Lemma of Definition~\ref{def:ultrafilter-lemma}, and
let \(I\) be infinite.
\begin{enumerate}
  \item The cofinite subsets of \(I\) form the Fr\'echet filter.
  \item Some ultrafilter on \(I\) extends that filter, and every such
        ultrafilter is nonprincipal.
  \item An ultrafilter on \(I\) is nonprincipal exactly when it contains
        every cofinite subset of \(I\).
\end{enumerate}
In particular, \(\N\) carries a nonprincipal ultrafilter.
\index{Frechet filter@Fr\'echet filter}
\end{hblemma}

\section{Ultraproducts and \L o\'s's method}
\label{sec:p03-ultraproducts}

For this section assume the Axiom of Choice in the ambient metatheory.  It
makes products of nonempty domains nonempty and supplies the coordinatewise
witness functions used at existential formulas.  It is not an axiom inside
any factor structure.

Let \((\mathcal M_i)_{i\in I}\) be \(\Lang\)-structures and let
\(\mathcal U\) be an ultrafilter on \(I\).  For
\(a,b\in\prod_{i\in I}\dom\mathcal M_i\), put
\[
 a\sim_{\mathcal U}b
 \quad\Longleftrightarrow\quad
 \{i\in I:a(i)=b(i)\}\in\mathcal U.
\]

\begin{hbdefinition}{def:p03-ultraproduct}{Ultraproduct}
The \emph{ultraproduct} is written
\[
 \prod_{i\in I}\mathcal M_i/\mathcal U.
\]
Its domain is the set of \(\sim_{\mathcal U}\)-classes
\([a]_{\mathcal U}\).  Constants
and functions are interpreted coordinatewise.  For every \(m\)-ary relation
symbol \(R\),
\[
 R([a_1]_{\mathcal U},\ldots,[a_m]_{\mathcal U})
 \quad\Longleftrightarrow\quad
 \{i:\mathcal M_i\entails
 R(a_1(i),\ldots,a_m(i))\}\in\mathcal U.
\]
If all factors are \(\mathcal M\), this structure is the \emph{ultrapower}
\(\mathcal M^I/\mathcal U\).  The quotient is meaningful because the next
lemma proves independence of representatives.
\index{ultraproduct}
\index{ultrapower}
\end{hbdefinition}

\begin{hblemma}{thm:p03-ultraproduct-well-defined}{proved=annex:p3.types.4}{Ultraproduct well-definedness}
The relation \(\sim_{\mathcal U}\) is an equivalence relation, and every
coordinatewise constant, function, and relation interpretation in
Definition~\ref{def:p03-ultraproduct} is independent of the chosen
representatives.
\end{hblemma}

\begin{hblemma}{thm:p03-los-terms}{proved=annex:p3.types.5}{Evaluation of terms}
Let \(t(x_1,\ldots,x_n)\) be a term.  For representatives
\(a_1,\ldots,a_n\), define \(b_t(i)=t^{\mathcal M_i}
(a_1(i),\ldots,a_n(i))\).  Then
\[
 t^{\prod_i\mathcal M_i/\mathcal U}
 ([a_1]_{\mathcal U},\ldots,[a_n]_{\mathcal U})=[b_t]_{\mathcal U}.
\]
\end{hblemma}

\begin{hbmethod}{def:p03-los-method}{The \L o\'s induction}
For a formula \(\theta(\bar x)\) and representatives \(\bar a\), track the
truth set
\(X_\theta(\bar a)=\{i:\mathcal M_i\entails\theta(\bar a(i))\}\).
First use term evaluation to reduce equality and relation atoms to the
quotient definitions.  Next translate each Boolean connective into the
corresponding complement, intersection, or union of truth sets and apply the
ultrafilter alternative.  For \(\exists y\,\psi\), an ultraproduct witness
\([b]_{\mathcal U}\) makes \(X_\psi(b,\bar a)\) large and hence makes the
existential truth set large.  Conversely, from a large existential truth set,
Choice selects a witness in each satisfying factor and an arbitrary element
elsewhere; the resulting \([b]_{\mathcal U}\) is an ultraproduct witness.
Universal formulas may then be treated as negated existentials.  These are
the load-bearing cases of the full proof in Annex~\ref{annex:p3.types.6}.
\index{Los induction@\L o\'s induction}
\end{hbmethod}

\begin{hbtheorem}{thm:p03-los}{proved=annex:p3.types.6}{Los's theorem@\L o\'s's theorem}
Assume the Axiom of Choice as above.  For every first-order
\(\Lang\)-formula \(\varphi(x_1,\ldots,x_n)\) and representatives
\(a_1,\ldots,a_n\),
\[
 \prod_{i\in I}\mathcal M_i/\mathcal U
 \entails\varphi([a_1]_{\mathcal U},\ldots,[a_n]_{\mathcal U})
\]
if and only if
\[
 \{i\in I:\mathcal M_i\entails
 \varphi(a_1(i),\ldots,a_n(i))\}\in\mathcal U.
\]
\end{hbtheorem}

\begin{hbcorollary}{thm:p03-diagonal-ultrapower}{proved=annex:p3.types.7}{Diagonal elementary embedding}
For every ultrapower \(\mathcal M^I/\mathcal U\), the map
\(d:\dom\mathcal M\to\dom(\mathcal M^I/\mathcal U)\) sending \(a\) to the
class of the constant function with value \(a\) is an elementary embedding.
\end{hbcorollary}

If \(\mathcal M_n\) is the finite linear order of size \(n\) and
\(\mathcal U\) is nonprincipal on \(\N\setminus\{0\}\), the ultraproduct is
infinite: for each \(k\), the sentence asserting at least \(k\) elements holds
on a cofinite, hence \(\mathcal U\)-large, set of indices.

\begin{hbwarning}{def:p03-los-scope}{First-order scope of \L o\'s's theorem}
The theorem transfers first-order formulas only.  In particular, ``the domain
is finite'' is not first-order definable over all structures, so an infinite
ultraproduct of finite structures is no contradiction.
\end{hbwarning}

\begin{hbtheorem}{thm:p03-keisler-shelah}{cited=Shelah1971Ultrapowers}{Keisler--Shelah isomorphism theorem}
For structures \(\mathcal M\) and \(\mathcal N\) in the same signature,
\(\mathcal M\equiv\mathcal N\) if and only if some ultrapower of
\(\mathcal M\) is isomorphic to some ultrapower of \(\mathcal N\).
No fixed index set, ultrafilter, nonprincipality condition, or cardinal bound
is asserted.  The reverse implication follows from \L o\'s's theorem and
isomorphism invariance; the forward implication is Shelah's theorem.
\end{hbtheorem}

\begin{hbtheorem}{thm:p03-compactness-ultraproduct}{proved=annex:p3.types.8}{Compactness via ultraproducts}
Assume the Axiom of Choice in the ambient set theory.  If every finite subset
of a set \(\Gamma\) of first-order \(\Lang\)-sentences is satisfiable, then
\(\Gamma\) is satisfiable.  The derivation in the annex uses the finite-subset
index set, extends its generated proper filter by the Ultrafilter Lemma, and
then applies \L o\'s's theorem; no choice-free claim is made.
\end{hbtheorem}

%% Migration accounting for v1.1 lines 6510--7149:
%% - every definition and every status-bearing result is retained above;
%% - the five explanatory examples are retained in compressed prose;
%% - the homogeneity remark and first-order warning are retained in substance;
%% - the five exercises are deliberately retired because they add practice
%%   prompts but no definitions, hypotheses, or result statements to this
%%   lookup chapter.
